% A2 on a 4d-continuum diamond waist. Two spacings. Not a -> 0 proof.
\documentclass[11pt]{article}
\usepackage[margin=1.05in]{geometry}
\usepackage{amsmath,amssymb,amsthm}
\usepackage{enumitem}
\usepackage[colorlinks=true,linkcolor=blue,citecolor=blue,urlcolor=blue]{hyperref}

\newcommand{\pP}{\textbf{[P]}}
\newtheorem*{admission}{Admission}

\title{\textbf{A2 on a $3{+}1$D Diamond Waist:\\
Two Lattice Spacings Do Not Refute Locality;\\
Bisognano--Wichmann Already Localizes the Continuum Wedge}}
\author{Robert W.\ McGwier\\
\small CTO, Cohere Technology Group\\
\small GitHub: \texttt{n4hy}}
\date{15 August 2026 --- Version 1}

\begin{document}
\maketitle

\begin{abstract}
\noindent
The same pre-registered locality gate as Papers~21--22:
$R(m)/R(0)<2$ for $0<mL\le 8$. The field is a free staggered
fermion in three-space: the equal-time slice of a $3{+}1$D theory.
The region is the geodesic ball at the waist of a Minkowski causal
diamond, diameter $L=2R$. Local ansatz $=$ on-site $+$ six spatial
nearest neighbours.

Coarse $N=12$, $R=3$ ($123$ sites): ratios $1.36$, $1.33$.
Fine $N=16$, $R=5$ ($515$ sites): ratios $1.13$, $1.19$.
Auditor, different code, $N=14$, $R=4$: $1.19$, $1.22$.
C2 PASS on every spacing. \pP{} as a necessary-condition test.

This is not $a\to 0$, not the Standard Model, and not a selection
of de~Sitter. The continuum wedge theorem is already
Bisognano--Wichmann; the CFT-ball theorem is already
Casini--Huerta--Myers. Both are cited, not re-derived.
\end{abstract}

\section{Why $3{+}1$D, not a $4$-space lattice}

Papers~21--22 tested $d=1,2,3$ spatial lattices. The missing object
named there was a \emph{$4$d continuum diamond}, not a
four-dimensional spatial grid. The virtual modes of a $3{+}1$D
field live on the Minkowski diamond. Their equal-time restriction
is a ball in three-space. A $4$-spatial lattice would be a
different theory.

\section{Gate}

Unchanged. $K=\log((1-C_A)/C_A)$ from the occupied projector of
\[
H_{ii}=m(-1)^{x+y+z},\qquad H_{\langle ij\rangle}=-1.
\]
Anything beyond the lattice adjacency cannot be a spacetime scalar
$X$ in Jacobson's split. Two spacings at matched $mL\in\{3,6\}$
are the continuum \emph{check}, locked before the run.

\section{Numbers}

\begin{center}
\begin{tabular}{lcccc}
\hline
spacing & $n_{\mathrm{ball}}$ & $R(0)$ & $R/R(0)$ at $mL=3$ & at $mL=6$ \\
\hline
$N=12$, $R=3$ & $123$ & $0.132$ & $1.36$ & $1.33$ \\
$N=14$, $R=4$ (audit) & --- & $0.150$ & $1.19$ & $1.22$ \\
$N=16$, $R=5$ & $515$ & $0.158$ & $1.13$ & $1.19$ \\
\hline
\end{tabular}
\end{center}

The ratio stays under $2$. It is smaller on the finer grid
(\emph{computed}, two points). That is the direction a continuum
limit would need. It is not the limit.

An independent adversary, different grids $N=13,15$, denser $mL$
scan, no shared code: worst ratio $1.29$ (at $mL=4$). Cube-shaped
waist also stayed under $2$. Catch, not a fail: $R(m)$ is not
monotone. It peaks near $mL\sim 4$--$5$ and then falls. Mass does
not monotonically localize $K$ on this stencil.

\section{What the continuum already says}

Bisognano--Wichmann \cite{BW1975}: the modular Hamiltonian of a
Rindler wedge in Minkowski, including $3{+}1$D, is
\[
K=2\pi\int_{x^1>0}x^1\,T_{00}\,d^3x.
\]
Casini--Huerta--Myers \cite{CHM2011}: for a ball in a CFT,
\[
K=2\pi\int_{|x|<R}\frac{R^2-r^2}{2R}\,T_{00}\,d^3x.
\]
Both are local. The lattice test is the massive deformation, which
those theorems do not cover.

\section{What this does not do}

It does not take $a\to 0$. It does not install the Standard Model.
It does not fix $\eta=1/4G$. It does not copy FGHMV onto de~Sitter
(Paper~17). It does not promote Jacobson 2016 to \pP{} (Paper~18).
Einstein$+\Lambda>0$ is already de~Sitter as a metric vacuum
(Paper~19); entanglement still does not select the sign of $\Lambda$.

\begin{admission}
I pushed A2 onto the $3{+}1$D diamond waist I had named as the
next object. I did not write that a two-spacing lattice is the
continuum, and I did not write that entanglement gravity is
de~Sitter.
\end{admission}

\begin{thebibliography}{9}
\bibitem{BW1975}
J.J.~Bisognano and E.H.~Wichmann,
J.\ Math.\ Phys.\ 16 (1975) 985.

\bibitem{CHM2011}
H.~Casini, M.~Huerta and R.C.~Myers,
JHEP 05 (2011) 036,
arXiv:1102.0440.

\bibitem{Jacobson2016}
T.~Jacobson,
Phys.\ Rev.\ Lett.\ 116 (2016) 201101,
arXiv:1505.04753.
\end{thebibliography}

\end{document}
